\chapter{Épreuve chronométrée 7}\label{ch:iv1:timed-test-7}

\begin{ivtrains}
Il faut la traiter d'une seule traite, sans notes ni calculatrice.
\end{ivtrains}

\ivtestsection{A}{6}{40}

\begin{question}\label{iq:iv1:timed-test-7:1}
Vous lancez sans cesse un dé équilibré. En moyenne, combien de lancers faut-il pour obtenir à la fois un 1 et un 2~? Et quelle est cette espérance jusqu'à ce que vous ayez vu les trois faces paires~?
\end{question}

\begin{question}\label{iq:iv1:timed-test-7:2}
Une pièce dont la probabilité de faire pile vaut $0.1$ est lancée 10 fois. Quelle est la probabilité d'obtenir un nombre pair de piles~? Vous pouvez utiliser $0.8^{10} \approx 0.107$.
\end{question}

\begin{question}\label{iq:iv1:timed-test-7:3}
$X$ et $Y$ ont chacun une variance égale à 1, et leur corrélation vaut $\frac12$. Quelle est la variance de $X + 2Y$~? Quelle est la pente de la régression de $Y$ sur $X + Y$~?
\end{question}

\begin{question}\label{iq:iv1:timed-test-7:4}
Deux dés équilibrés affichent $X$ et $Y$. Quelle est l'espérance de $\max(X - Y, 0)$~?
\end{question}

\begin{question}\label{iq:iv1:timed-test-7:5}
Deux dés équilibrés sont lancés. Sachant qu'au moins l'un d'eux affiche un 6, quelle est la probabilité que les deux affichent un 6~? Si vous savez en revanche que le premier dé affiche un 6, quelle est la probabilité que les deux affichent un 6~?
\end{question}

\begin{question}\label{iq:iv1:timed-test-7:6}
Une matrice carrée d'ordre 2 possède quatre coefficients, chacun valant indépendamment 0 ou 1 avec la même probabilité. Quelle est la probabilité qu'elle soit inversible sur les réels~?
\end{question}

\ivtestsection{B}{6}{60}

\begin{question}\label{iq:iv1:timed-test-7:7}
Un sac contient trois pièces~: une pièce équilibrée, une pièce à deux piles, et une pièce dont la probabilité de pile vaut $\frac34$. Vous en tirez une au hasard et la lancez deux fois, en obtenant deux piles. Quelle est la probabilité que le lancer suivant de cette même pièce donne pile~?
\end{question}

\begin{question}\label{iq:iv1:timed-test-7:8}
La durée de vie d'une ampoule suit une loi uniforme sur $[0, 2]$ années. Chaque ampoule est remplacée dès qu'elle tombe en panne, et ce processus dure depuis très longtemps. À un instant quelconque, quelle est l'espérance de l'âge de l'ampoule en service~? Quelle est l'espérance de la durée de vie totale de cette ampoule~?
\end{question}

\begin{question}\label{iq:iv1:timed-test-7:9}
$Z$ suit une loi normale centrée réduite. En utilisant $\varphi(1) = 0.242$ et $P(Z > 1) = 0.159$, déterminez $\mathbb{E}[Z \mid Z > 1]$. Une distribution des pertes suit une loi normale de moyenne 0 et de volatilité 1. Sachant que le quantile à 99~\% vaut $2.326$ et que $\varphi(2.326) = 0.0267$, quelle est l'expected shortfall à 99~\%~?
\end{question}

\begin{question}\label{iq:iv1:timed-test-7:10}
Une permutation de $\{1, 2, 3, 4\}$ est choisie au hasard. Une inversion est un couple de positions $i < j$ dont les valeurs sont mal ordonnées. Déterminez le nombre moyen d'inversions, ainsi que sa variance.
\end{question}

\begin{question}\label{iq:iv1:timed-test-7:11}
$X$ et $Y$ sont des variables normales centrées réduites de corrélation $\frac12$. Quelle est $\mathbb{E}\max(X, Y)$~?
\end{question}

\begin{question}\label{iq:iv1:timed-test-7:12}
Vous lancez 100 dés équilibrés. Estimez la probabilité que la somme soit au moins égale à 380, en admettant qu'un dé ait pour variance $\frac{35}{12}$ et que $P(Z > 1.73) = 0.042$ pour une variable $Z$ normale centrée réduite.
\end{question}
